
\subsubsection{6.1 Benchmark Compression as a Measurable Phenomenon}

The result that 31.1\% of 466 qualifying benchmarks exhibit at least one compression event over 7 years provides the first systematic cross-domain prevalence estimate for this phenomenon. The Granger causality result (p = 1.854 × 10⁻⁵ at lag 2) establishes temporal precedence of submission accumulation over score variance compression; the reverse direction is not confirmed. Together these results support the interpretation that submission accumulation drives compression rather than the reverse, though Granger causality does not rule out confounding variables such as benchmark age or domain adoption dynamics.

The dissociation between the weak Spearman correlation (ρ = 0.052, significant at p = 1.51 × 10⁻⁵ but far below the ρ > 0.4 pre-registered target) and the strong Granger signal is consistent with a nonlinear threshold mechanism: compression occurs once submission count crosses a benchmark-specific threshold, producing negligible contemporaneous correlation but detectable lagged temporal structure.

\subsubsection{6.2 H_d Signal Discriminability Caveat}

The H-E1 results show large effect sizes (|d| > 5) in all three domains. These were obtained on synthetic panels calibrated to real data statistics, not on real benchmark data with real saturation labels. Effect sizes of this magnitude on synthetic data may be optimistic: idealized synthetic panels lack real-world confounds including publication recency bias, task heterogeneity, submission selectivity, and noise in saturation labels. Real-data validation (proposed as FW1 in the research plan) would re-run the H-E1 signal computation on the H-M1 real panel using \texttt{compression_event} labels as the saturation proxy. The real-data effect sizes are expected to be positive but cannot be assumed to exceed |d| = 5.

\subsubsection{6.3 Collapse Event Operationalization as a Methodological Finding}

The H-M2 gate failure provides a methodological finding of independent value: expanding Kendall τ on quarterly score variance is structurally incompatible with the compression signal established in H-M1. Any research group using expanding rank correlation on the same column as the compression signal will encounter the same near-zero event count. This root-cause analysis prevents investment in an operationalization that is guaranteed to fail on any quarterly leaderboard panel where compression is defined as variance reduction.

The proposed resolution (FW2) is to redefine collapse as \texttt{compression_event == 1.0} for at least 4 consecutive quarters. Based on the 389 events across 145 benchmarks in H-M1, this criterion is expected to yield approximately 100 collapse events — sufficient to train a Cox PH model and conduct Kaplan-Meier lead-time analysis.

\subsubsection{6.4 Limitations}

\textbf{L1 — Synthetic data in H-E1 (Severity: Moderate).} Discriminability validation used synthetic panels because the PWC REST API was unavailable. The |d| > 5 results are on idealized data. H-M1 independently confirms the compression mechanism on real data, but H_d signal performance on real saturation labels has not been measured.

\textbf{L2 — Cox PH survival model not validated (Severity: High).} The central prospective prediction component — C-index ≥ 0.70 under time-split validation and lowest-quintile hazard ratio ≥ 2× — was not evaluated. H-M3 and H-M4 were blocked by the H-M2 operationalization failure. This paper does not provide a validated benchmark health prediction system.

\textbf{L3 — Temporal ordering of H_d signals not established (Severity: High).} The H-M2 gate failure means temporal precedence of H_d signals relative to discriminative collapse has not been established. The cross-sectional NLP AUC = 0.857 confirms that H_d^NLP can identify currently-compressed benchmarks but does not establish that the signal precedes compression.

\textbf{L4 — Granger panel power (Severity: Moderate).} The minimum time series length filter (≥ 9 quarters) reduced testable benchmarks from 466 to 41. The 12.2\% individual significance rate is likely an underestimate of the true mechanism prevalence. The panel-level Granger result (p = 1.854 × 10⁻⁵) passes Bonferroni correction for 41 tests and is the primary Granger finding.

\textbf{L5 — CV H_d direction not normalized (Severity: Low-Moderate).} The directional asymmetry between CV (lower = saturated) and NLP/tabular (higher = saturated) means current H_d signals cannot be directly combined into a multi-domain scalar. Per-domain sign normalization is required for the Cox model covariate design.

\textbf{L6 — Scope limited to CV and NLP from PWC (Severity: Moderate).} The qualifying benchmark filter (≥ 20 submissions, ≥ 2 years) yielded 466 benchmarks from the PWC archive covering CV and NLP. Tabular domain results in H-E1 come from synthetic panels; tabular real-data analysis was not performed. Results are not validated for RL, biomedical NLP, or other leaderboard platforms (Kaggle, HELM, BIG-Bench).

\subsubsection{6.5 Future Work}

\textbf{Collapse event recalibration (FW2, immediate priority).} Redefining collapse as compression_event ≥ 4 consecutive quarters would yield approximately 100 events from the existing panel, enabling H-M3 execution without new data collection.

\textbf{Real-data H_d validation (FW1).} Re-running H-E1 signal computation on the H-M1 real panel with compression_event as the saturation label would establish whether the synthetic |d| > 5 results are representative of real-data discriminability.

\textbf{Per-domain sign normalization (FW4).} Implementing direction normalization (negate H_d^CV before combining) is required for the multi-domain Cox model covariate.

\textbf{Cox PH model with nonlinear cumulative_count encoding (FW3).} The threshold mechanism suggested by the Spearman-Granger dissociation implies that \texttt{cumulative_count} should be encoded as a threshold-spline or piecewise-linear predictor rather than a linear one in the Cox model.

\textbf{Extension to additional domains and platforms (FW7).} The domain phenotype framework (convergence vs. divergence saturation) provides a conceptual basis for designing H_d signals for RL benchmarks and other leaderboard platforms.

